% Unit 36 — Named constants and one-line substitutions
\input{technique-note.tex}
\title{Integration Bee --- Unit 36:\\ Named Constants and One-Line Substitutions}
\author{UK Uni Integration Bee prep\\ A parameter, a substitution, or a named constant}
\date{}
\begin{document}
\maketitle
\section*{Status}
Tags: \texttt{yang:6.9}; techniques: \texttt{unique-tricks}.
Yang \S6.9 is a drawer of leftovers. In this drill set every integral is one of three things: a parameter integral from unit 34, a substitution $u=\ln x$, or Catalan's constant left as Catalan's constant. There is no fourth trick to memorize, and there is no elementary closed form hiding behind $G$. The route is \texttt{docs/techniques/concepts/habits.md}. The catalog id \texttt{catalan} stays unticked.
\section{Official anchors}
\begin{itemize}
\item Unit 34 for $\int_0^1 (x^b-x^a)/\ln x\,dx=\ln((b+1)/(a+1))$.
\item Unit 35 for $\int_0^{\pi/2}\ln(\sin x)\,dx=-(\pi/2)\ln 2$.
\item Catalan's constant $G=\sum_{k=0}^{\infty}(-1)^k/(2k+1)^2$. Syllabus \S2 lists it. Yang Topics Guide 2nd Ed \S6.9 Unique Tricks.
\end{itemize}
\section{Learning goals}
\begin{enumerate}
  \item Recognise $(x^b-x^a)/\ln x$ and integrate the geometric parameter. Do not series-expand $\ln x$ in the denominator.
  \item Leave $G$ as $G$. Do not replace it by a rational number or a logarithm.
  \item For $1/(x\ln x)$, substitute $u=\ln x$.
\end{enumerate}

\section{Worked parameter --- \texttt{U36-D01}}
\[
\int_0^1\frac{x-1}{\ln x}\,dx=\ln 2.
\]
For $a>-1$, $\int_0^1 x^a\,dx=1/(a+1)$. Integrate in $a$ from 0 to 1:
\[
\int_0^1\frac{x-1}{\ln x}\,dx=\ln\frac{1+1}{0+1}=\ln 2.
\]
\texttt{U36-D07} runs the parameter from 0 to 2 and stores $\ln 3$. That is \texttt{U34-D04}. The improper endpoint is the same limit used in unit 34.

\section{Catalan's constant --- \texttt{U36-D03}}
\[
\int_0^1\frac{\arctan x}{x}\,dx=G.
\]
On $[0,1)$, $\arctan x=\sum_{k=0}^{\infty}(-1)^k x^{2k+1}/(2k+1)$, so
\[
\frac{\arctan x}{x}=\sum_{k=0}^{\infty}(-1)^k\frac{x^{2k}}{2k+1}.
\]
Integrate termwise. Each term contributes $(-1)^k/(2k+1)^2$. The series at $x=1$ is the definition of $G$. The integrand extends continuously through $x=0$ by the value $1$. The stored answer is the constant $G$, written \texttt{Catalan} in the drill feed. Stop there.

\section{The logarithm substitution --- \texttt{U36-D04}}
\[
\int\frac{dx}{x\ln x}=\ln|\ln x|+C.
\]
Let $u=\ln x$, so $du=dx/x$. \texttt{U36-D08} is the partner $\int(\ln x)/x\,dx$. The same $u$ gives $\int u\,du=\frac12(\ln x)^2$, the stored antiderivative.

\section{Faded definite logarithms --- \texttt{U36-D05} and \texttt{U36-D06}}
On $[2,4]$,
\[
\bigl[\ln|\ln x|\bigr]_2^4=\ln(\ln 4)-\ln(\ln 2)=\ln(2\ln 2)-\ln(\ln 2)=\ln 2,
\]
because $\ln 4=2\ln 2$. That is the stored value of \texttt{U36-D05}.

On $[e,e^e]$,
\[
\bigl[\ln|\ln x|\bigr]_e^{e^e}=\ln\bigl(\ln(e^e)\bigr)-\ln(\ln e)=\ln e-\ln 1=1.
\]
That is the stored value of \texttt{U36-D06}.

\section{Faded parameter --- \texttt{U36-D07}}
\[
\int_0^1\frac{x^2-1}{\ln x}\,dx.
\]
For $a>-1$, $\int_0^1 x^a\,dx=1/(a+1)$. Integrate in $a$ from $0$ to $2$:
\[
\ln(2+1)-\ln(0+1)=\ln 3.
\]
That is the stored value, and it is \texttt{U34-D04}. The same formula with the upper limit $4$ is a scale, not a drill:
\[
\int_0^1\frac{x^4-1}{\ln x}\,dx=\ln 5.
\]
SymPy checks $\ln 5$. The parameter from $0$ to $3$ is already stored as \texttt{U11-D15}, value $\ln 4$. Leave $G$ as $G$. The fence $\ln 5$ is a parameter value.

\section{Already taught --- \texttt{U36-D02}}
\[
\int_0^{\pi/2}\ln(\sin x)\,dx=-\frac\pi2\ln 2
\]
is \texttt{U35-D06}. Use the double-angle argument in that note. Calling it a unique trick does not change the proof.

\section{Drill pointers}
See \texttt{drills.tex} / \texttt{drills.md}. Feed: \texttt{verify/drills\_unit36.json}.

\section*{On a miss}
Teach this file. If the answer is $G$, say so and stop. If the denominator is $\ln x$, build the parameter or the substitution $u=\ln x$. Do not invent an answer that is not in \texttt{drills.md}.
\end{document}
